% Source: https://artofproblemsolving.com/community/c3h2836699_aime_2021_q8_help_need
% Stable identity: aops-topic-2836699-q01
Let $g(x)=|20|x|-x^2|$. The equation says $g(x)=c+21$ or $g(x)=c-21$.

The graph of $g$ is even. For $x\ge0$, $20x-x^2$ rises from $0$ to $100$ at $x=10$ and falls to $0$ at $x=20$, then is negative; taking the absolute value reflects the part beyond $x=20$ upward. So for $0<k<100$ the line $y=k$ meets the graph for $x>0$ in $3$ points, hence $g(x)=k$ has $6$ real solutions; for $k>100$ it has $2$; for $k=100$ it has $4$; and for $k=0$ it has $3$ ($x=0,\pm20$).

To get $12$ solutions we need both $c+21$ and $c-21$ in $(0,100)$ (they are distinct, so the solution sets are disjoint): $21<c<79$. That gives $c=22,\dots,78$, i.e. $\boxed{57}$ integers.
